% Six-page version. Saved predictions/metrics were audited on 2026-11-05.
% AP uses maxDets=300. The image-disjoint split is NOT video-group-disjoint.
% Provenance and unresolved records: BENCHMARK_EVIDENCE_AUDIT.md.
% !TeX program = xelatex
% Build: latexmk -xelatex -interaction=nonstopmode PhysicalAI_For_Optimal_Feeding_Control_6p.tex
\documentclass[10pt,conference,letterpaper]{IEEEtran}
\usepackage{newtxmath}
\usepackage{fontspec}
\setmainfont[
  SmallCapsFont=texgyretermes-regular.otf,
  SmallCapsFeatures={Letters=SmallCaps}
]{Times New Roman}
\setmonofont{Courier New}
\usepackage{geometry}
\geometry{letterpaper,top=18mm,bottom=18mm,left=16.5mm,right=16.5mm}
\setlength{\columnsep}{5mm}
\usepackage{graphicx}
\graphicspath{{assets/}{Physical-AI-For-Optimal-Feeding-Control/assets/}{../figures/}{figures/}}
\usepackage{float}
\usepackage{amsmath}
\usepackage{tikz}
\usetikzlibrary{arrows.meta,positioning,calc}
\usepackage{tabularx}
\usepackage{booktabs}
\usepackage{cite}
\usepackage{microtype}
\usepackage{balance}
\usepackage{url}
\pagestyle{empty}
\raggedbottom
\setlength{\intextsep}{4pt plus 2pt minus 2pt}
\setlength{\abovecaptionskip}{2pt}

\title{A Physical AI Framework for Adaptive Feeding Based on Visual Feed Dynamics}
\author{\IEEEauthorblockN{\fontsize{11}{13}\selectfont Sin-Ze Yang, Cheng-Hsun Chuang, Yong-Xiang Chen\textsuperscript{*}}
\IEEEauthorblockA{\fontsize{10}{12}\selectfont\itshape Undergraduate Program in Intelligent Computing and Big Data\\
Chung Yuan Christian University, Taiwan\\
{\fontsize{9}{11}\selectfont\normalfont\ttfamily *Contact: yxchen@cycu.edu.tw}}}

\begin{document}
\maketitle
\begin{abstract}
Adaptive aquaculture feeding requires feedback on environmental responses to dispensing actions. This study proposes a Physical AI framework that integrates object detection, causal Kalman estimation, and embedded actuation to guide subsequent portions based on visible-feed dynamics. A bounded decision policy uses estimated changes in visible-feed area after dispensing. To evaluate the perception module, six YOLO26 and RF-DETR configurations are compared on 99 test images. RF-DETR-m achieves the highest overall AP50--95 of 0.363, while YOLO26-m achieves the highest feed F1 of 0.66. YOLO26-s achieves a feed F1 of 0.64 with a recorded median local pipeline latency of 5.00 ms. These results highlight the importance of feed-specific performance in perception model selection. However, all models fall below the provisional feed recall target of 0.75, and overlapping video groups across data splits limit generalization. Temporal signal reliability, deployment timing, and feeding effectiveness remain unvalidated.
\end{abstract}

\section{Introduction}
Excess feed increases aquaculture costs and organic loading. Fixed schedules cannot respond to changes in appetite associated with temperature, dissolved oxygen, fish size, or health. A feedback-based feeder therefore needs observations that can inform the next dispensing action.

Beyond static observations, changes in visible feed after dispensing provide a physical response that can inform subsequent feeding decisions. Using this response requires reliable temporal measurements and a validated relationship between disappearance and feeding demand; detector accuracy alone cannot establish either.

Here, \emph{Physical AI} denotes the proposed perception--action loop in which dispensing changes the environment observed for the next decision. The framework connects visual feed dynamics, causal temporal estimation, and bounded actuation. The completed perception benchmark is presented separately from the proposed control design.

The main contributions are:
\begin{enumerate}
\setlength{\itemsep}{0pt}
\setlength{\parsep}{0pt}
\setlength{\topsep}{0pt}
\item A Physical AI framework that connects visual feed dynamics with subsequent dispensing decisions.
\item A causal temporal estimator and bounded feeding policy with observation quality and safety constraints.
\item An empirical comparison of six detection configurations to assess feed-specific accuracy and computational efficiency for the proposed perception module.
\end{enumerate}

\section{Related Work}
Vision-based feeding research uses residual pellets or fish activity as feedback. Foster \emph{et al.} detect and count uneaten pellets~\cite{foster1995}; Liu \emph{et al.} measure feeding activity in recirculating systems~\cite{liu2014}; Zhou \emph{et al.} combine near-infrared vision with feeding decisions~\cite{zhou2018}. Our proposed signal is the relative change in visible feed after a portion, whose interpretation still requires controlled validation.

RF-DETR combines a DINOv2-style backbone with neural architecture search~\cite{robinson2026}, while YOLO26 supports native NMS-free inference~\cite{jocher2026}. We compare their native pipelines for dense feed detection and distinguish feed-specific performance from overall multiclass accuracy. Temporal estimation is addressed separately using a standard causal Kalman filter~\cite{kalman1960}.

\section{System Architecture}
\label{sec:architecture}
\subsection{Hardware and Software Overview}
In the reference design, a Raspberry Pi Zero 2 W captures and streams images of a fixed surface region. The server runs RF-DETR inference, image-quality checks, causal Kalman filtering, and feeding decisions. The Pi relays commands to a BMduino that drives the servo dispenser (Fig.~\ref{fig:hardware}). RF-DETR is the reference detector, not a final deployment selection. Figure~\ref{fig:device} shows the prototype; integrated operation has not been verified by deployment traces.

\begin{figure}[H]
\centering
\includegraphics[width=.48\columnwidth]{feeding_device.jpg}
\caption{Prototype feeder with BMduino, Raspberry Pi, and servo-driven outlet.}
\label{fig:device}
\end{figure}

\begin{figure*}[!t]
\centering
\includegraphics[width=\textwidth]{smart_feeding_architecture.pdf}
\caption{Reference architecture. The server runs RF-DETR inference, image-quality checks, causal Kalman filtering, and feeding decisions; the Pi captures/streams images and relays commands; the BMduino controls the servo. Solid arrows indicate data exchange; dashed and dash-dotted arrows indicate control requests and hardware control, respectively.}
\label{fig:hardware}
\end{figure*}

% Keep subsequent figures in the float queue after the wide architecture figure.
\subsection{Control State Machine}
The proposed server-side controller cycles through \textsc{Feed} (request a portion), \textsc{Wait} (exclude motor activity and falling feed), \textsc{Observe} (collect $L\geq2$ scheduled samples), and \textsc{Decide} (apply Section~\ref{sec:decision}). Camera/network failure, prolonged invalid observations, or a missing actuator response requires \textsc{Suspend}. A local Pi command timeout must also inhibit further dispensing when the server is unreachable. These failure responses are requirements awaiting tests.

\subsection{Timing of One Feeding Cycle}
The planned logs associate frame/cycle IDs with capture times, detections, quality flags, state transitions, commands, and replies. The Pi monotonic clock provides the reference timeline; cross-device events require clock mapping with reported uncertainty or Pi send/receive timestamps. Command acknowledgements alone do not establish physical motor motion.

Video calibration must determine the waiting interval. During \textsc{Observe}, $L$ scheduled samples at rate $f_s$ span $(L-1)/f_s$ seconds. Missing or rejected samples retain their scheduled positions. The signal update and decision eligibility rules below apply to this fixed-duration window.

\section{Signal Model and Estimation}
The proposed estimator tracks visible-feed area over time to provide a temporal signal for subsequent feeding decisions.

\subsection{Visible-Feed Measurement}
For sample $t$, let $N(t)$ be the number of detections, with box $B_i$, width $w_i$, height $h_i$, confidence $c_i$, and class $y_i$. Boxes are expressed on the fixed camera-image pixel grid after reversing model-input resizing. For a fixed region of interest $\Omega$, the visible-feed proxy is
\begin{equation}
\begin{aligned}
A_t={}&\sum_{i=1}^{N(t)}w_i h_i\mathbb{I}(y_i=\mathrm{feed})\\
&\cdot\mathbb{I}(c_i\geq\tau)\mathbb{I}(B_i\cap\Omega\neq\varnothing).
\end{aligned}
\label{eq:area}
\end{equation}
$A_t$ sums whole-box areas for accepted feed detections intersecting $\Omega$; it does not clip boxes to the ROI or remove their overlap. It is an area proxy, not pellet mass or an exact count. Figure~\ref{fig:detection} illustrates class selection only.

\begin{figure}[!tbp]
\centering
\includegraphics[width=.50\columnwidth]{detection.png}
\caption{Example overlay illustrating feed-class selection. Equation~\eqref{eq:area} counts whole feed boxes intersecting the ROI; this image is not an accuracy evaluation.}
\label{fig:detection}
\end{figure}

\subsection{Quality Gate}
The proposed quality gate considers feed confidence, occlusion, glare, and filter innovation. Rejected or missing frames receive prediction-only updates and do not count as accepted observations. Gate features and thresholds require target-system calibration, including tests that distinguish an empty scene from missed feed detections. The control threshold $\tau$ is also uncalibrated; the benchmark confidence threshold does not establish a suitable control setting.

\subsection{Causal Kalman Filter}
Let $x_t$ represent latent visible-feed area and $z_t=A_t$ its measurement. The scalar random-walk Kalman model~\cite{kalman1960} is
\begin{equation}
\begin{aligned}
x_t&=x_{t-1}+q_t,\quad q_t\sim\mathcal{N}(0,Q),\\
z_t&=x_t+v_t,\quad v_t\sim\mathcal{N}(0,R).
\end{aligned}
\label{eq:process}
\end{equation}
\begin{equation}
\begin{aligned}
\hat x_t^-&=\hat x_{t-1},\quad P_t^-=P_{t-1}+Q,\\
K_t&=P_t^-/(P_t^-+R).
\end{aligned}
\label{eq:prediction}
\end{equation}
\begin{equation}
\begin{aligned}
\hat x_t&=\hat x_t^-+K_t(z_t-\hat x_t^-),\\
P_t&=(1-K_t)P_t^-,\quad \hat A_t=\max(0,\hat x_t).
\end{aligned}
\label{eq:posterior}
\end{equation}
Superscript $-$ denotes prediction, and $P_t$ is the state-error variance. Equations~\eqref{eq:prediction}--\eqref{eq:posterior} give the update for an accepted observation. For a rejected or missing sample, $\hat x_t=\hat x_t^-$ and $P_t=P_t^-$. The nonnegative output $\hat A_t$ is then used by the decision rule. Variances $Q,R$ and initialization remain to be calibrated; reducing $Q/R$ trades responsiveness for smoothing.

Figure~\ref{fig:feed_trace} compares estimates on a supplied RF-DETR sequence. Its offline Rauch--Tung--Striebel (RTS) smoother uses future observations and is excluded from control. Without manual reference areas and error/delay measurements, the trace illustrates processing behavior rather than estimator accuracy.

\begin{figure*}[!t]
\centering
\includegraphics[width=.96\textwidth]{feed_per_frame.png}
\caption{Illustrative RF-DETR feed-area sequence with causal Kalman filtering, offline RTS smoothing, and an exponential moving average (EMA, $\alpha=0.1$). Area is in pixels$^2$; the horizontal axis is frame number. No manual reference is available for this trace.}
\label{fig:feed_trace}
\end{figure*}

\section{Decision Rule}
\label{sec:decision}
The proposed policy adjusts subsequent feeding portions according to the estimated visible-feed disappearance after dispensing, subject to observation validity and safety constraints.

Let $u_k$ be the dispensing command for cycle $k$ and $g(u_k)$ its mass under a calibration that remains to be established. After \textsc{Wait}, let $t$ index the last of the $L$ scheduled observation samples. If either endpoint is invalid or the accepted-sample count is insufficient, the controller enters \textsc{Suspend}. Otherwise, it computes the relative disappearance ratio
\begin{equation}
\rho_t=\frac{\hat A_{t-L+1}-\hat A_t}{\max(\hat A_{t-L+1},\varepsilon)},\qquad L\in\mathbb{N},\ L\geq2.
\label{eq:rho}
\end{equation}
Here, $\varepsilon>0$ is an area floor that limits division by near-zero estimates; it does not replace the initial-area eligibility check. A larger positive $\rho_t$ means a greater relative decrease over the fixed window, not a per-second rate or a physiological hunger measure.

With $\theta_L<\theta_H$, $\rho_t\geq\theta_H$ qualifies for an increase, $\theta_L\leq\rho_t<\theta_H$ selects a nominal baseline portion, and $\rho_t<\theta_L$ calls for reduction or suspension. The baseline, thresholds, minimum accepted-sample count, and low-initial-area cutoff remain calibration choices. An initial area below that cutoff cannot qualify for an increase. Frames from \textsc{Feed}/\textsc{Wait} do not enter this calculation.

\subsection{Portion Update and Safety Limits}
For active feeding, the proposed bounds are $g_{\min}\le g(u_{k+1})\le g_{\max}$, together with a daily mass cap and a maximum increase per cycle. An increase requires two consecutive eligible high-disappearance windows; reductions need no such persistence. \textsc{Suspend} issues no dispensing command, so active-portion bounds do not require a nonzero portion in that state. These rules constrain commands; their feeding effectiveness remains untested.

\section{Perception Module Evaluation}
\label{sec:compare}
The benchmark compares YOLO26-n/s/m (Y-n/s/m) and RF-DETR-n/s/m (R-n/s/m) using archived records and metrics recomputed from saved predictions. Training, inference, and GPU timing were not rerun. Table~\ref{tab:training} reports training settings.

\subsection{Dataset}
The dataset combines 199 images from \texttt{fishtracker\_v2} version 1~\cite{fishtracker_dataset} and 313 from \texttt{feeding\_fish} version 2~\cite{feedingfish_dataset}. Source export metadata records auto-orientation and stretch resizing to $640\times640$ and $512\times512$, respectively, without export-time augmentation. All 512 source-image hashes match the manifest. Table~\ref{tab:labels} gives the verified label mapping; only unified \texttt{feed} detections (ID 1) enter Eq.~\eqref{eq:area}.

% Primary snapshot: audit/benchmark_evidence/ (captured read-only from MyLinux).
% Source labels, configs, split/evaluation code, prediction files, and run records
% are retained with capture_receipt.json; audit/recomputed_metrics.json records
% in-memory recomputation without modifying the remote benchmark.
\begin{table}[H]
\caption{\textsc{Unified Class Identifiers and Source-Label Mapping}}
\label{tab:labels}
\centering
\footnotesize
\setlength{\tabcolsep}{3pt}
\begin{tabularx}{\columnwidth}{@{}c l >{\raggedright\arraybackslash}X@{}}
\toprule
\textbf{ID} & \textbf{Label} & \textbf{Source dataset: original label} \\ \midrule
0 & \texttt{debris} & \texttt{fishtracker\_v2}: \texttt{debris} \\
1 & \texttt{feed} & \texttt{fishtracker\_v2}: \texttt{feed}; \texttt{feeding\_fish}: \texttt{pellet} \\
2 & \texttt{fish} & \texttt{fishtracker\_v2}: \texttt{fish}; \texttt{feeding\_fish}: \texttt{feeding-fish} \\
3 & \texttt{lumpfish} & \texttt{fishtracker\_v2}: \texttt{lumpfish} \\ \bottomrule
\end{tabularx}
\end{table}

The seed-42 manifest assigns 314 images to training and 99 each to validation and testing. Training contains 123 images from the first source and 191 from the second; validation and test each contain 38 and 61, respectively. Recounting source annotations confirms 24,467 boxes (Table~\ref{tab:dataset}).

Image hashes and exact decoded pixels are disjoint across splits. However, images are allocated \emph{within} source/video strata, and all eight filename-derived groups span all splits. We therefore describe this as a same-video image holdout; temporal similarity remains a leakage risk.

\begin{table}[H]
\caption{\textsc{Bounding-Box Counts Before Training Augmentation}}
\label{tab:dataset}
\centering
\footnotesize
\setlength{\tabcolsep}{3pt}
\begin{tabularx}{\columnwidth}{@{}l*{4}{>{\centering\arraybackslash}X}@{}}
\toprule
\textbf{Label} & \shortstack{\textbf{Train}\\(314 images)} & \shortstack{\textbf{Val}\\(99 images)} & \shortstack{\textbf{Test}\\(99 images)} & \textbf{Total} \\ \midrule
\texttt{debris} & 305 & 102 & 70 & 477 \\
\texttt{feed} & 11,748 & 2,839 & 3,380 & 17,967 \\
\texttt{fish} & 3,602 & 1,178 & 1,156 & 5,936 \\
\texttt{lumpfish} & 52 & 17 & 18 & 87 \\ \midrule
\textbf{Total} & \textbf{15,707} & \textbf{4,136} & \textbf{4,624} & \textbf{24,467} \\ \bottomrule
\end{tabularx}
\end{table}

Table~\ref{tab:class_images} counts images containing each class. Classes may co-occur, so counts cannot be summed to obtain split sizes. Recounting confirms one training image with no annotated boxes. Both tables precede training augmentation.

\begin{table}[H]
\caption{\textsc{Images Containing Each Class Before Augmentation}}
\label{tab:class_images}
\centering
\footnotesize
\begin{tabularx}{\columnwidth}{@{}l*{3}{>{\centering\arraybackslash}X}@{}}
\toprule
\textbf{Label} & \textbf{Train} & \textbf{Val} & \textbf{Test} \\ \midrule
\texttt{debris} & 49 & 12 & 13 \\
\texttt{feed} & 160 & 45 & 49 \\
\texttt{fish} & 263 & 85 & 82 \\
\texttt{lumpfish} & 36 & 10 & 13 \\ \bottomrule
\end{tabularx}
\end{table}

\begin{table*}[!t]
\caption{\textsc{Training Settings Verified from Run Records and Logs}}
\label{tab:training}
\centering
\footnotesize
\setlength{\tabcolsep}{4pt}
\renewcommand{\arraystretch}{1.12}
\begin{tabularx}{\textwidth}{@{}p{.21\textwidth} >{\raggedright\arraybackslash}X >{\raggedright\arraybackslash}X@{}}
\toprule
\textbf{Setting} & \textbf{YOLO26 (n/s/m)} & \textbf{RF-DETR (n/s/m)} \\ \midrule
Input size & 640/640/640 & 384/512/576 \\
Epoch cap / patience & 100 / 20 & 100 / 20; minimum improvement 0.001 \\
Logged epochs completed & 100/100/100 & 49/88/97 (zero-based maxima + 1) \\
Micro-batch $\times$ accumulation & $16\times1$ / $16\times1$ / $8\times2$; effective batch 16 & $1\times16$ for all sizes; effective batch 16 \\
Optimizer & SGD; LR 0.01, momentum 0.937, weight decay $5\times10^{-4}$ & AdamW; base LR $10^{-4}$, encoder LR $1.5\times10^{-4}$, weight decay $10^{-4}$ \\
LR schedule & Non-cosine; final LR factor 0.01; warm-up 3 epochs & Step schedule, drop epoch 100; no warm-up \\
Selected checkpoint & \texttt{best.pt}, native validation fitness & \texttt{checkpoint\_best\_total.pth}, best native validation AP across regular/EMA weights \\ \bottomrule
\end{tabularx}
\end{table*}

\subsection{Training and Preprocessing}
The export code converts polygons to enclosing axis-aligned boxes and remaps classes using Table~\ref{tab:labels}. YOLO uses normalized box coordinates and unified class IDs; the COCO evaluation export uses pixel coordinates and category IDs offset by one. Both trainers use the same manifest.

Training and split seeds are 42. YOLO checkpoints record Ultralytics 8.4.53; the audited environment contains RF-DETR 1.7.0 and PyTorch 2.12.0. Source snapshots and checkpoint hashes are retained because the recorded Git commit omits benchmark code.



Saved YOLO arguments specify horizontal/vertical flips 0.5/0.2, rotation $\pm15^{\circ}$, translation 0.15, scale 0.6, HSV 0.02/0.8/0.5, mosaic 1.0 (disabled for the final 10 epochs), MixUp 0.2, and perspective 0.001. RF-DETR records multiscale/expanded-scale training without a custom augmentation override; the captured package default enables horizontal flipping at 0.5.

Checkpoint selection uses each trainer's native validation criterion. RF-DETR validation has a 500-detection limit, distinct from the common benchmark protocol below. Together with different resolutions, augmentations, and completed epochs, this makes the experiment a pipeline comparison rather than an isolated test of architecture.

\subsection{Comparison Protocol}
All six run records share the manifest hash. Their saved predictions cover the same 99 test image IDs and are evaluated against 4,624 ground-truth boxes. Recomputing the metrics reproduces every stored detection value and confidence interval.

Prediction adapters use score floor 0.001 and at most 300 detections per image. AP is computed with pycocotools 2.0.11 over IoU 0.50:0.05:0.95 and \texttt{maxDets=[1,10,300]}, extracting the 300-detection column. Overall AP averages the four classes. All AP values in this paper use that limit rather than the standard COCO limit of 100.

Precision, recall, and F1 use confidence 0.25 and class-aware greedy one-to-one matching at IoU 0.5, pooling counts across images. Percentile 95\% intervals resample images 2,000 times with seed 123. They quantify image-sampling variation conditional on this split and omit video dependence and training-seed variation.
% Primary values: audit/benchmark_evidence/results/benchmark/runs/<run>/metrics.json.
% Recomputed AP@300 and AP@100: audit/recomputed_metrics.json.
% Legacy local summary AP50--95 values match recomputed AP@100; all paper AP
% values and plots now use AP@300 consistently. P/R/F1 use confidence 0.25.

\subsection{Detection Accuracy and Efficiency}
\begin{table*}[t]
\caption{\textsc{Detection Accuracy on the 99-Image Same-Video Test Split}}
\label{tab:compare}
\centering
\footnotesize
\setlength{\tabcolsep}{3pt}
\renewcommand{\arraystretch}{1.15}
\begin{tabularx}{\textwidth}{@{}l c*{8}{>{\centering\arraybackslash}X}@{}}
\toprule
\textbf{Run} & \textbf{Input} & \shortstack{\textbf{Overall}\\\textbf{AP50--95}} & \shortstack{\textbf{Overall}\\\textbf{AP50}} & \shortstack{\textbf{Feed}\\\textbf{AP50--95}} & \shortstack{\textbf{Feed}\\\textbf{Precision}} & \shortstack{\textbf{Feed}\\\textbf{Recall}} & \shortstack{\textbf{Feed}\\\textbf{F1}} & \shortstack{\textbf{Fish}\\\textbf{AP50--95}} & \shortstack{\textbf{Fish}\\\textbf{F1}} \\ \midrule
Y-n & 640 & 0.219 & 0.456 & 0.192 & \shortstack{0.66\\{[0.62, 0.69]}} & \shortstack{0.52\\{[0.42, 0.59]}} & \shortstack{0.58\\{[0.50, 0.63]}} & 0.258 & \shortstack{0.54\\{[0.50, 0.58]}} \\
Y-s & 640 & 0.271 & 0.578 & 0.234 & \shortstack{0.72\\{[0.67, 0.75]}} & \shortstack{0.58\\{[0.47, 0.66]}} & \shortstack{0.64\\{[0.56, 0.69]}} & 0.284 & \shortstack{0.61\\{[0.56, 0.64]}} \\
Y-m & 640 & 0.307 & 0.581 & 0.258 & \shortstack{0.69\\{[0.65, 0.72]}} & \shortstack{0.64\\{[0.53, 0.71]}} & \shortstack{0.66\\{[0.59, 0.71]}} & 0.303 & \shortstack{0.61\\{[0.56, 0.65]}} \\
R-n & 384 & 0.277 & 0.453 & 0.086 & \shortstack{0.46\\{[0.40, 0.51]}} & \shortstack{0.31\\{[0.26, 0.37]}} & \shortstack{0.37\\{[0.32, 0.42]}} & 0.323 & \shortstack{0.61\\{[0.57, 0.64]}} \\
R-s & 512 & 0.358 & 0.578 & 0.160 & \shortstack{0.54\\{[0.47, 0.60]}} & \shortstack{0.51\\{[0.45, 0.58]}} & \shortstack{0.53\\{[0.47, 0.57]}} & 0.357 & \shortstack{0.63\\{[0.58, 0.67]}} \\
R-m & 576 & 0.363 & 0.597 & 0.184 & \shortstack{0.55\\{[0.48, 0.61]}} & \shortstack{0.57\\{[0.50, 0.63]}} & \shortstack{0.56\\{[0.50, 0.61]}} & 0.369 & \shortstack{0.65\\{[0.61, 0.69]}} \\ \bottomrule
\end{tabularx}
\par\smallskip
\parbox{\textwidth}{\footnotesize Y: YOLO26; R: RF-DETR. All AP values use maxDets=300. Brackets are image-level bootstrap 95\% intervals (2,000 resamples, seed 123); AP intervals are absent. Video groups overlap across train/validation/test, so these intervals do not establish unseen-video performance.}
\end{table*}

Table~\ref{tab:compare} and Fig.~\ref{fig:accuracy} show different rankings for overall and feed-specific accuracy. RF-DETR-m has the highest overall AP50--95 (0.363), whereas YOLO26-m has the highest feed AP50--95 (0.258), recall (0.64), and F1 (0.66). These are descriptive single-seed differences; no paired tests establish statistical significance.

\begin{figure*}[!t]
\centering
\includegraphics[width=.98\textwidth]{accuracy_overview.png}
\caption{Detector accuracy on the same-video test split: (a) overall AP50--95, (b) feed AP50--95, and (c) feed F1 at confidence 0.25. AP uses maxDets=300. Blue: YOLO26; orange: RF-DETR. F1 error bars are image-bootstrap 95\% intervals; AP intervals are unavailable.}
\label{fig:accuracy}
\end{figure*}

\begin{table*}[!t]
\caption{\textsc{Recorded Local Pipeline Latency and GPU Memory}}
\label{tab:efficiency}
\centering
\footnotesize
\setlength{\tabcolsep}{3pt}
\renewcommand{\arraystretch}{1.1}
\begin{tabularx}{\textwidth}{@{}l*{7}{>{\centering\arraybackslash}X}@{}}
\toprule
\textbf{Run} & \shortstack{\textbf{Model-only}\\\textbf{p50 (ms)}} & \shortstack{\textbf{Model-only}\\\textbf{p95 (ms)}} & \shortstack{\textbf{Local pipeline}\\\textbf{p50 (ms)}} & \shortstack{\textbf{Local pipeline}\\\textbf{p95 (ms)}} & \shortstack{\textbf{Reciprocal p50}\\\textbf{(s$^{-1}$)}} & \shortstack{\textbf{Peak allocated}\\\textbf{GPU (MiB)}} & \textbf{Input} \\ \midrule
Y-n & 3.47 & 4.02 & 4.60 & 5.54 & 218 & 70 & 640 \\
Y-s & 3.85 & 4.29 & 5.00 & 5.91 & 200 & 126 & 640 \\
Y-m & 8.22 & 8.34 & 9.45 & 9.77 & 106 & 197 & 640 \\
R-n & 6.43 & 7.05 & 10.35 & 11.18 & 97 & 176 & 384 \\
R-s & 10.76 & 11.58 & 14.74 & 15.62 & 68 & 204 & 512 \\
R-m & 12.58 & 13.40 & 16.65 & 17.81 & 60 & 222 & 576 \\ \bottomrule
\end{tabularx}
\par\smallskip
\parbox{\textwidth}{\footnotesize RTX 5060 Laptop GPU; inference batch 1, FP32 model-only tensors; 50 warm-up and 200 timed iterations on one image. Reciprocal p50 is a derived rate, not measured streaming throughput. Local pipeline excludes file I/O, network, and actuation. Retained timing values were restored from CSV, not independently remeasured.}
\end{table*}

Table~\ref{tab:efficiency} contains historical timing values restored from CSV. The retained profiler synchronizes CUDA after each iteration, uses the first test image for all models, and measures peak allocated GPU memory in MiB. Full historical profile blocks and clock-state logs are absent, so these values remain provisional.

Local-pipeline timing includes native preprocessing and postprocessing, starting from BGR arrays for YOLO and preconverted RGB PIL images for RF-DETR. The RF-DETR input conversion is outside the timer. YOLO uses fused convolution/batch-normalization in PyTorch without TensorRT; RF-DETR has no inference optimization. These boundaries differ from complete acquisition-to-actuation timing.

\begin{figure*}[!t]
\centering
\includegraphics[width=.90\textwidth]{accuracy_vs_latency.png}
\caption{Accuracy versus recorded median local-pipeline latency: (a) feed F1 with image-bootstrap 95\% intervals; (b) overall AP50--95 at maxDets=300. Test video groups overlap training. Local timing excludes acquisition, file I/O, network, actuation, and RF-DETR input conversion.}
\label{fig:pareto}
\end{figure*}

\subsection{Implications for Adaptive Feeding}
The perception module supplies visible-feed measurements for temporal estimation and feeding decisions. The results show that overall detection accuracy does not necessarily reflect feed-specific performance, suggesting that model selection should prioritize task-relevant accuracy and computational efficiency. YOLO26-s offers a potential accuracy--latency tradeoff for further evaluation (Fig.~\ref{fig:pareto}). However, all configurations remain below the provisional feed recall target of 0.75. Further validation is needed to determine how detection errors affect visible-feed estimation and subsequent feeding decisions.

\section{Discussion and Validation Priorities}
\label{sec:validation}
\subsection{Design Implications for Physical AI Feeding}
The proposed framework and perception benchmark suggest three design considerations.

\subsubsection{Task-Oriented Perception}
Perception models should be evaluated according to downstream decision requirements. The observed differences between overall and feed-specific accuracy highlight the importance of task-relevant detection performance.

\subsubsection{Temporal Observation Separation}
Separating dispensing from observation through the proposed state machine may reduce interference from falling feed and motor activity. Quality gating further limits the influence of unreliable measurements.

\subsubsection{Safety-Constrained Adaptation}
Feeding decisions should incorporate observation validity, portion limits, and suspension conditions to prevent unreliable measurements or system failures from triggering uncontrolled dispensing.

These considerations inform the proposed system design but require validation through integrated feeding experiments.

\subsection{Signal and Deployment Validation}
Group-disjoint retraining and testing on new videos are needed: the same-video split and image-bootstrap intervals do not establish generalization or training-seed stability.

Visible-feed disappearance can reflect sinking, transport outside $\Omega$, glare, occlusion, or box overlap. No-fish controls and independently observed feeding events must distinguish these effects from consumption.

Calibration on sessions excluded from final evaluation must fix camera geometry, dispensing mass $g(u_k)$ through repeated weighings, filter initialization and $Q,R$, gate settings, $\tau,f_s,L$, and all portion-policy parameters.

Compare raw $A_t$, EMA, and causal Kalman estimates against manual references over identical windows using area error, delay, and incorrect feeding decisions; lower variance alone is insufficient. Restrict offline smoothing to analysis.

Deployment tests must measure acquisition-to-actuation delay, sampling rate, power, and the failure responses in Section~\ref{sec:architecture}. Physical feedback must confirm dispensing; acknowledgements and repeated-image GPU timing are insufficient.

\subsection{Tank-Level Outcomes}
Planned experiments will randomize replicate tanks with matched initial biomass to timer-based or adaptive feeding, measuring feed mass, visible uneaten-feed duration, biomass gain, mortality, water quality, and feed conversion ratio (FCR), with effect sizes and confidence intervals. Feeding and environmental benefits remain unvalidated.

\section{Conclusion}
This study proposes a Physical AI framework that uses visual feed dynamics to inform adaptive feeding decisions. The perception benchmark demonstrates differences between overall and feed-specific detection performance, supporting task-oriented model evaluation. While the framework provides a reference design for integrating perception, temporal estimation, and physical actuation, further validation of signal reliability, system operation, and feeding outcomes is required.

\begin{thebibliography}{8}
\bibitem{foster1995} M. Foster, R. Petrell, M. R. Ito, and R. Ward, ``Detection and counting of uneaten food pellets in a sea cage using image analysis,'' \emph{Aquacultural Engineering}, vol. 14, no. 3, pp. 251--269, 1995.
\bibitem{liu2014} Z. Liu, X. Li, L. Fan, H. Lu, L. Liu, and Y. Liu, ``Measuring feeding activity of fish in RAS using computer vision,'' \emph{Aquacultural Engineering}, vol. 60, pp. 20--27, 2014.
\bibitem{zhou2018} C. Zhou \emph{et al.}, ``Near infrared computer vision and neuro-fuzzy model-based feeding decision system for fish in aquaculture,'' \emph{Computers and Electronics in Agriculture}, vol. 146, pp. 114--124, 2018.
\bibitem{robinson2026} I. Robinson, P. Robicheaux, M. Popov, D. Ramanan, and N. Peri, ``RF-DETR: Neural Architecture Search for Real-Time Detection Transformers,'' in \emph{Proc. ICLR}, 2026.
\bibitem{jocher2026} G. Jocher \emph{et al.}, ``Ultralytics YOLO26: Unified Real-Time End-to-End Vision Models,'' arXiv:2606.03748, 2026. [Online]. Available: \url{https://arxiv.org/abs/2606.03748}
\bibitem{kalman1960} R. E. Kalman, ``A new approach to linear filtering and prediction problems,'' \emph{Journal of Basic Engineering}, vol. 82, no. 1, pp. 35--45, 1960.
\bibitem{fishtracker_dataset} JTElectric, ``FishTracker\_v2,'' \emph{Roboflow Universe}, ver. 1, exported May 17, 2026. [Online]. Available: \url{https://universe.roboflow.com/jtelectric/fishtracker_v2/dataset/1}
\bibitem{feedingfish_dataset} ``feeding\_fish,'' \emph{Roboflow Universe}, ver. 2, exported May 19, 2026. [Online]. Available: \url{https://universe.roboflow.com/aifeeder-70gjh/feeding_fish/dataset/2}
\end{thebibliography}
\end{document}
